\chapter{Scope and actual checklist delta}
\label{component:p0be1d6088cc6}
\noindent\textbf{Component }\path{p0be1d6088cc6}\par
{\small\noindent Source: \path{BASE/proofs/routing/routing_audit.tex}\par}
{\small\noindent SHA256: \path{0be1d6088cc6370c7cb2ec4b2b0ed2a294c691a0cf96af071b5f14d9a8e9250e}\par}
{\small\noindent Inventory use: C1 input and finite dispatch; C1.4 subdivision | C14\_\allowbreak{}ENTIRE\_\allowbreak{}ORIGINAL\_\allowbreak{}PARENT\_\allowbreak{}GRAPH; precise subordinate uses retained below\par}
\medskip
\section*{Scope and actual checklist delta}

This note closes the two technical audit imports left in U1.A6 by the
R18 routing review: constrained curvature with one or two hard events,
and the four-matching terminal with arbitrary internal coefficient
scales. It also writes an explicit common input, the source-entry maps,
and an exhaustive decision tree into the historical thirteen rows and
their C14/R1--R10 subdivisions. All original physicality, contact,
coefficient signs, strict gaps, literal zeros, and total-scale conditions
are carried along as conjuncts.

The distinction between a proved dispatch implication and an accepted
terminal certificate is essential. The complete positive-total-scale
terminal bundle in A5 is only partially repaired: the recovered finite-edge
proof and the new \path{a5_recovery/finite_pitchfork_fit.tex} cover the
ordinary/finite-pitchfork subbundle with $d,C,A,D,\alpha$ units and its
exact symmetry images. The companion
\path{a5_recovery/weak_alpha_unit_n_fit.tex} supplies a second precise
C2 subdomain: $d,n,m,C,A,D$ all units, $\alpha$ weak, and
$\lambda=\beta=0$. The rest of weak-$\alpha$, projective-pitchfork,
and reciprocal terminals remain imported. This audit also
found a specific mismatch in the preserved uniform C1 theorem statements:
the printed infinitesimal-$e$ and unbounded-$e$ theorems do not by
themselves cover unit $e$. The companion proof
\path{us_unit/bounded_delta_all_tests.tex} repairs that subentry directly.
It extends the exact full-parent contact--conductor proof to bounded
$\Delta=1/e$ and bounded testing $w$, proves positive coefficient weights
without comparing independent rates, and supplies the complete testing
join. Thus an old all-$e$ assertion is no longer needed for that subentry.

The exact parent and Laurent determinant identities and the separated
bounded strict-Gram lemma remain the already-admitted inputs. The
focused routing checker has 225 passing exact assertions, covering both full-parent
reflections, support
classification, the two-hard activation table, finite source-type
routing, and four-matching normal completion. The separate bounded-$\Delta$
checker has 33 passing exact conductor and full-source lift assertions.
The scoped positive-scale repair has a recovered 15-identity, 30-entry
full-parent replay and 11 new exact projective-boundary checks.
Its C2 weak-$\alpha$ subdomain has a separate recovered 9-identity,
31-entry full-tail replay and 4 exact state-jet assertions.
These assertions do not
prove an unresolved commuting-boundary terminal. No global tensor-SOS
or spectrahedral-shadow conclusion follows from this note.

All cited paths are given in the accompanying \path{ROUTING_LEDGER.json},
relative to the recovered \path{checkpoint_R19} directory, with
SHA-256 hashes; new continuation files have their own explicit path root.
The symbols \texttt{R13}, \texttt{OLD}, and \texttt{INH}
in the ledger have one explicit expansion each; no absent external
manuscript is substituted for a preserved file.

\section{One exact input, before optional exclusions}

Let $K_c,K_t\supset\R$ be independent real closed fields. The bounded
subring of either field consists of elements whose absolute value is
at most some real number; its maximal ideal consists of infinitesimals.
For $q>0$ write $q:\I$ when $q$ is infinitesimal, $q:\U$ when $q$
and $q^{-1}$ are bounded, and $q:\Otype$ when $q^{-1}$ is infinitesimal.
These are exhaustive, disjoint types. Bounded means $\I$ or $\U$;
noninfinitesimal means $\U$ or $\Otype$. Zero is always a separate
literal case. Type notation is external shorthand for ordered-field
inequalities relative to real bounds, not a new algebraic equation.

In the coefficient order $(d,\lambda,n,m,C,A,D,\alpha,\beta)$ and
with $D_u=1+u^2$, $D_v=1+v^2$, retain the complete original parent
\begin{align}
F_c={}&4\lambda D_u(W+vX)^2+4nD_uD_vX^2
 +4m(vW-X+uvY-uZ)^2\notag\\
&+4CD_v(uW-Y)^2+4AD_uD_vY^2
 +(4D+2d)(uW+uvX-Y-vZ)^2\notag\\
&+(4\alpha+2d)D_u(vY-Z)^2
 +(4\beta+2d)D_v(X+uZ)^2-2dD_uD_vZ^2. \tag{P}
\end{align}
Its symmetric state matrix is $P_c(u,v)$, so $F_c=\psi^TP_c\psi$.
Physicality is the exact quantified polynomial condition
$P_c(u,v)\succeq0$ for every source-field projective pair, interpreted
in the four homogeneous angle charts. It is never replaced by
positivity of a truncated scalar or a modified parent.

The canonical remaining source input $Q(c)$ is the following conjunction:
\begin{enumerate}
\item $c\in K_c^9$, $d>0$, all eight event coefficients nonnegative,
and physicality of (P).
\item $\delta_{\rm tot}=d/(d+\lambda+n+m+C+A+D+\alpha+\beta):\I$.
This quotient is computed before any later $d=1$ or $C=2$ normalization.
\item The source is coercive and has five or six positive events.
There is a finite positive contact $u_c,v_c>0$ and nonzero $\psi_c$
such that $P_c\psi_c=0$ and
$\partial_u F_c=\partial_v F_c=0$ at this contact.
\item A deterministic omission convention has chosen C1 or C2 as
described below. Its literal coefficient equalities and positivity
conditions are part of $Q$, not conditions to be inferred later from
a limiting approximation.
\end{enumerate}
The testing input is an arbitrary independent $K_t$ projective angle
pair and an arbitrary real or complex state. A complex state is split
into real and imaginary parts because the parent is real symmetric.
No source contact is identified with a testing point.

If one wants precisely the input after A6, append that at least one
testing measurement pair commutes in real residue. We prove the source
dispatch with all tests allowed, so that additional restriction cannot
create a missing source case. Optional historical exits may likewise be
intersected afterwards: if $Q\Rightarrow E\lor\bigvee R_i$, then
$Q\land\neg E\land B\Rightarrow\bigvee(R_i\land B)$ for every
common additional predicate $B$. This elementary implication avoids
using an unspecified phrase such as ``the complement of all prior
admissions'' as the premise of the coverage proof itself.

\paragraph{Canonical omission and symmetry convention.}
Use the fixed cycle
$(\lambda,D,C,A,\alpha,m,\beta,n)$.
The A3 reduction gives compatible omissions: pairs have cycle distance
one or two, and triples consist of three consecutive vertices.
For a triple, choose a distance-one pair first, with the lexicographically
least pair in this fixed cycle breaking ties. Choose the first word in
a fixed enumeration of the finite Bell relabeling group that sends it
to $\{\lambda,n\}$. The extra zero is then $D=0$ or $\beta=0$.
For two omissions, distance one enters C1 and distance two enters C2.
Thus the two canonical domains are
\begin{align*}
\mathrm{C1}:&\quad \lambda=n=0,\quad m,C,A,\alpha>0,\quad
 D,\beta\ge0,\quad D+\beta>0;\\
\mathrm{C2}:&\quad \lambda=\beta=0,\quad m,C,A,D,\alpha>0,
 \quad n\ge0.
\end{align*}
The five-active C2 boundary $n=0$ is already assigned C1 by the
precedence convention, but it is deliberately allowed in C2 theorem
statements and all later implications. No proof is lost by this overlap.
The partition is disjoint for chosen canonical records, not for their
unoriented physical-source images.

\section{A6: the formerly imported technical audit}

\subsection{Weighted centers do not require comparable hard weights}

At a noncommuting testing residue, use an entangled residue state as an
invertible two-by-two matrix. A fixed local ray has a unique orthogonal
partner on the other side. Consequently its vanishing losing events
form a matching of the eight-cycle, of size at most four. At most three
such product-event rows are independent state coordinates. Product
residue states instead have strictly positive signed-core value at a
noncommuting pair; their contribution is a pure-source positive factor
times a pure-testing square, plus event squares.

Suppose independent hard event amplitudes are exact coordinates $x_i$,
their positive weights are $a_i$, and the next coefficient scale is
$\sigma$, with $\varepsilon_i=\sigma/a_i:\I$. The local polynomial is
\[
F=\sum_i a_i b_i(y)x_i^2+\sigma G(x,y),\qquad b_i(y_0)>0.
\]
Set $x_i=\varepsilon_i s_i$ in the coefficient-field stationary
equations. After division by $\sigma$ their normal equations have
residue $2b_i(y)s_i+G_{x_i}(0,y)=0$, and their tangent equations have
residue $\nabla_yG(0,y)=0$. If the constrained soft Hessian is positive
definite, the residue Jacobian is invertible: its hard diagonal entries
are $2b_i>0$ and its soft Schur block is precisely that Hessian.
The semialgebraic implicit-function statement over $\R$, followed by
transfer, supplies an exact source-field stationary center. Physicality
makes its value nonnegative.

The exact finite Taylor roots are
$(x_i-x_{ic})/\sqrt{\varepsilon_i}$ and $y-y_c$.
For each hard packet, $x_{ic}=\varepsilon_i s_{ic}$ makes its
center-square and center-linear contributions bounded after this
scaling. All remaining higher-degree factors are actual deviations,
not quotients by a source opening. The residue Gram is strict, so the
admitted separated strict-Gram lemma applies. None of these operations
uses a mixed inverse or orders a source infinitesimal against a testing
infinitesimal.

\subsection{Exact deficient-support classification}

For $d>0$ write the eight outer Laurent coefficients of the determinant
as $q_e\ge0$, $e\in\{-1,0,1\}^2\setminus\{0\}$. At a stationary
zero its logarithmic Hessian is $\sum_e w_e ee^T$ with $w_e>0$ for
active exponents. It fails to be positive definite only if the support
has rank zero or rank one. A stationary nonempty rank-one support must
contain opposite exponents; one-sided support cannot have zero gradient.

For one hard event, take $\lambda$ by symmetry and extract
$\partial q_e/\partial\lambda$. For two hard matching events the three
orbits have representatives $\{\lambda,m\}$, $\{\lambda,C\}$,
and $\{\lambda,\alpha\}$. Extract the coefficient of their product.
A state basis adapted to the independent hard amplitudes shows that this
leading coefficient is a positive angular factor times the constrained
state determinant. The positivity table below is therefore exact.
Juxtaposition means both coefficients are positive.
\[
\begin{array}{c|ccc}
e&\{\lambda,m\}&\{\lambda,C\}&\{\lambda,\alpha\}\\\hline
(-1,-1)&\alpha&\alpha&\text{always}\\
(-1,0)&D\text{ or }\alpha&D\text{ or }\alpha m&D\text{ or }m\\
(-1,1)&D&Dm&Dm\\
(0,-1)&n\text{ or }A\alpha&n&A\text{ or }n\\
(0,1)&\beta\text{ or }DC&D\text{ or }\beta m&DC\text{ or }\beta m\\
(1,-1)&An&0&An\\
(1,0)&AC\text{ or }\beta n&A&AC\text{ or }\beta n\\
(1,1)&\beta C&\beta&\beta C
\end{array}
\]
Inspection of the four opposite-pair lines gives no rank-one stationary
support for $\{\lambda,m\}$. For the other two pairs the only
possibility is the perfect matching
$\{\lambda,C,\alpha,\beta\}$. The one-hard calculation has exactly
that same possibility. With no hard event there are precisely the two
perfect matchings. These are positivity-support implications and allow
every omitted weight to be literally zero.

All rank-zero supports are contained in a nonextendable matching triple,
whose cyclic gaps are $(2,3,3)$. For one hard $\lambda$, representatives
are $\{\lambda,A,\beta\}$, $\{\lambda,m,C\}$, and
$\{\lambda,m,A\}$. For the pair $\{\lambda,m\}$ only the latter
two occur; for $\{\lambda,C\}$ only $\{\lambda,m,C\}$ occurs;
for $\{\lambda,\alpha\}$ rank zero is impossible. At $u=v=1$, the
states $(1,-1,0,2)$ and $(1,-1,1,3)$ annihilate respectively the
$\{\lambda,m,A\}$ and $\{\lambda,m,C\}$ events and give signed-core
values $-4$ and $-8$. Cycle symmetry supplies the remaining triples.
Hence a physical constrained source with $d>0$ cannot have rank zero.

On a perfect-matching curve all matching event squares and the signed
core vanish. For the representative $\{n,m,A,D\}$ the four independent
event-normal differentials satisfy
\[
dH=2t\{(1+t^2)dX-(1+t^2)dY-dR_m+dR_D\},\qquad t>0.
\]
Every coefficient is nonzero. Imposing at most two hard constraints
leaves a direction on which this derivative is nonzero; all event-square
derivatives vanish there. The constrained parent changes sign, contrary
to physicality. This removes the rank-one cases.

Thus every zero of the one/two-hard constrained determinant has positive
angular Hessian. The signed core is strictly positive on the $Z=0$
state subspace, so the constrained state matrix has respectively two
or one positive transverse state eigenvalues and a simple kernel with
$Z\ne0$. A Schur-complement calculation over the real residue field
identifies the angular Hessian with a positive multiple of the determinant
Hessian. The full soft Hessian is therefore strict. This is a real
residue calculation, not a Schur inverse in the tensor product.

For reproducibility the checker enumerates respectively 256, 128, and
64 supports for no hard event, one fixed hard event, and each two-hard
orbit. It checks every positivity-table entry on every support. The
enumeration corroborates the finite proof; it makes no numerical
positivity inference and does not test only generic positive coefficients.

\subsection{Four matching weights, all internal scales retained}

Suppose $\mu=\min(n,m,A,D)>0$, and $d/\mu$ and each of
$\lambda/\mu,C/\mu,\alpha/\mu,\beta/\mu$ are infinitesimal.
Rotate the cycle so $D=\mu$, and set
\[
\eta_1=\mu/n,\quad\eta_2=\mu/A,\quad\eta_3=\mu/m,
\qquad 0<\eta_i\le1.
\]
No positive lower bound is imposed on any $\eta_i$.
In the $W=1$ chart take $R=v-X+uvY-uZ$ as a state coordinate, so
$Z=(v-X+uvY-R)/u$. The fourth amplitude is $f+L$, where
\[
f=(u^2-v^2)/u,\quad
L=v(u+u^{-1})X-(1+v^2)Y+(v/u)R.
\]
At $X=Y=R=0$ the signed and fourth packets equal
\[
4\mu f^2+2duf
=4\mu\left(f+\frac{du}{4\mu}\right)^2-\frac{d^2u^2}{4\mu}.
\]
Testing $u=1$, $v=\sqrt{1+d/(4\mu)}$ proves, by physicality, that
$\tau=\max(\lambda,C,\alpha,\beta)>0$ and
$\rho=d^2/(\mu\tau)$ is bounded. Put $\varepsilon=d/\mu:\I$.

If $\rho:\I$, use the four normal roots
$\sqrt{n/\tau}X$, $\sqrt{A/\tau}Y$, $\sqrt{m/\tau}R$,
$\sqrt{\mu/\tau}(v-u)$, and the constant root. Their normal residue is
\[
4B^2z_1^2+4B^2z_2^2+4z_3^2+
4(-2z_4+B\sqrt{\bar\eta_1}z_1-B\sqrt{\bar\eta_2}z_2
             +\sqrt{\bar\eta_3}z_3)^2\succ0,
\quad B=1+u_0^2.
\]
The complementary packets give a strictly positive constant residue;
the signed linear corrections have coefficients $\sqrt\rho$ or
$\sqrt{\rho\eta_i}$ and vanish in residue.

If $\rho:\U$, substitute in source-field stationarity
\[
X=\varepsilon\eta_1s_1,\quad Y=\varepsilon\eta_2s_2,
\quad R=\varepsilon\eta_3s_3,\quad v-u=\varepsilon r,\quad u=t.
\]
With $B=1+t^2$, $\ell=(B,-B,1)$, $b=(B^2,B^2,1)$, and
$k=-2r+\sum\bar\eta_i\ell_i s_i$, the four divided normal equations
are
\[
8b_i s_i+8k\ell_i+(4tB,-4tB,0)_i=0,\qquad -16k-4t=0.
\]
Their Jacobian determinant is $32\cdot8^3B^4>0$, even if some
$\bar\eta_i=0$. Their unique solution is
\[
s_1=-\frac{t}{4B},\quad s_2=\frac{t}{4B},\quad s_3=\frac t4,
\quad r=\frac t8(1-\bar\eta_1-\bar\eta_2+\bar\eta_3).
\]
Let $P=\operatorname{res}((\lambda+\alpha)/\tau)$ and
$Q=\operatorname{res}((C+\beta)/\tau)$; $P+Q>0$. Exact normal
completion of the normalized residue objective leaves
\[
g(t)=4(1+t^2)(P+Qt^2)
 -\frac{\bar\rho t^2}{4}(1+\bar\eta_1+\bar\eta_2+\bar\eta_3).
\]
Physicality gives $g\ge0$. At a zero, stationarity of $g/(4t^2)$
gives $P/t^2=Qt^2$, forcing $P,Q>0$ and positive logarithmic curvature.
If $g(t_0)>0$, normal completion gives a strict constant-plus-normal
Gram. Otherwise adjoining the tangent equation gives an invertible
Jacobian and an exact stationary center with nonnegative value.

Put $\delta=v-u$. The full-parent Taylor roots are
\[
\sqrt{n/\tau}(X-X_c),\quad\sqrt{A/\tau}(Y-Y_c),\quad
\sqrt{m/\tau}(R-R_c),\quad\sqrt{\mu/\tau}(\delta-\delta_c),
\quad t-t_c.
\]
The packetwise bounds that justify the entire Taylor polynomial are
\[
\frac{a_i x_{ci}^2}{\tau}=\rho\eta_i s_{ci}^2,\qquad
\frac{a_i x_{ci}}{\sqrt{a_i\tau}}=\sqrt{\rho\eta_i}s_{ci}.
\]
The fourth normal has the same estimates with $\eta_i=1$. Complementary
and signed packets have no larger factors. Every packet is at most
quadratic in the state coordinates, so after two normal factors are
assigned the remaining higher factors are bounded deviations. The only
denominators are pure testing units $t$, together with admitted pure
coefficient units. The residue normal block is the strict form above;
its tangent Schur complement is $g''(t_0)/2>0$. This proves the full
packet restoration for arbitrary unequal matching weights.

At each earlier scale either a positive soft value or a strict constrained
Hessian finishes the argument, or another matching edge is added.
Three-edge kernels have the two inherited explicit Laurent forms; their
rank failures give respectively no stationary zero or a negative minimum
$2\sqrt{2D(d+2D)}-d-4D<0$ (squaring has difference $d^2>0$).
There can be at most four matching edges. The terminal above therefore
completes the finite matching-scale induction without a rate comparison.
This closes the specific A6 technical imports. It says nothing about a
commuting testing residue by continuity.

\section{Global sign entry and exact coordinate maps}

\subsection{C1}

Normalize $C=2$ by a positive pure-source factor. The finite contact has
the unique gauge
\[
(u_c,v_c;\psi_c)=(k,\ell;S,k(\theta-1),kSr,1),
\qquad S>0,\quad0<r,\theta<1.
\]
The inherited global sign proof has now had its curvature input audited
above. Over $\R$ it uses the unique maximal-$d$ contact graph over the
connected positive event orthant, a positive cofactor proving
$\theta>0$, and no-wall removal arguments for $Y=0$, $S-Y=0$, and
$\theta=1$. Its exact positive seed is
\[
(u,v;S,Y,\theta)=(1,1;1,1/4,1/2),\qquad
c=(27/2,0,0,4,2,9,5/4,13/4,9/4).
\]
The cross-multiplied amplitude identity extends from the nonzero Cramer
minor locus by analyticity on this specified graph. A boundary implicit
function argument opens $D=0$ or $\beta=0$, proves the strict signs
there, and only then transfers the finite first-order implication.
This use of real continuity proves source sign entry; it is not an
argument that tensor SOS is closed under non-Archimedean limits.

The full-parent Alice reflection is
\[
u'=1/u,\quad v'=v,\quad
G\psi=(vW-X,-W-vX,Z-vY,Y+vZ),
\]
with $(\lambda\ n)(m\ C)(A\ \alpha)(D\ \beta)$ on coefficients.
It obeys $G^TG=(1+v^2)I$ and
\[
u^2F_{c'}(1/u,v;G\psi)=(1+v^2)F_c(u,v;\psi).
\]
Its factors are pure testing squares after homogeneous normalization.
Apply it on the source side if needed to fix $0<k\le1$; equality is
kept in the unit case. Literal $D=0$ and $\beta=0$ exchange.

Define, exactly,
\[
\eta=\ell/(kS),\quad\zeta=S\ell/k,\quad e=k/S,
\quad x=k^2,\quad a=(1-r)/\eta,\quad t=\theta/\zeta,
\quad h=1-\theta.
\]
In particular $e=k\sqrt{\eta/\zeta}$, with the positive square root.
The original gaps are
\begin{gather*}
R=1+xr,\quad T=1+xh,\quad M=R-t>0,\quad J=T-a>0,\\
E=1-x\eta h>0,\quad Q=1-x\zeta r>0,\quad0<a,t<2.
\end{gather*}
Every branch below retains these exact relations. The same selected
minor remains nonzero when $D=0$ or $\beta=0$, so a boundary coefficient
column must not simply be deleted from the rational-source derivation.
The nonnegative source numerators $B_5,B_\beta$ in the ledger are
equivalent to the literal original inequalities $D,\beta\ge0$.

\subsection{C2}

Normalize $d=1$. At the source contact $Z=1$ the exact signs are
\[
W>0,\ X<0,\ Y>0,\ uW-Y>0,\ 1-vY>0,
\quad R_m>0,\ R_D<0,\ X+u>0,\ W+vX>0.
\]
The proof uses the analogous unique real maximal-contact graph, the
exact seed
\[
(n,m,C,A,D,\alpha)=(1/8,3/8,1/6,5/8,3/8,1/6),
\quad(u,v;\psi)=(1,1;1,-1/4,1/4,1),
\]
and event-amplitude removal. The wall $X=0$ would delete $n$ and enter
the C1 triple face, where the already-proved strict sign gives a
contradiction. The omitted-event signs follow from the original $X$
kernel equation and the exact C2 reflection. Literal $n=0$ is the C1
triple face directly. First-order transfer is performed after this real
argument.

The source variables and all strict gaps are
\begin{gather*}
c=u_c,\quad a=v_c/u_c,\quad b=aW_c,\quad r=-X_c/u_c,
\quad t=(u_cW_c-Y_c)/v_c,\\
h=aY_c/u_c=b-a^2t>0,\quad
g=R_m/u_c=(1+c^2)b-(1-r+a^2c^2t)>0,\\
j=-R_D/v_c=1+c^2r-t>0,\quad0<r<1,\quad
1-c^2h>0,\quad b-a^2c^2r>0.
\end{gather*}
The positive stationary determinant is a nonzero factor times
$ac^6rgt\,hj(1-c^2h)(1+c^2)^2(1+a^2c^2)^2\Delta$,
where
\[
\Delta=1+(1+z)b+z(1+z)b^2+2z(1+z)br+z^2r^2+2zr
 +a^2z\{(1+z)b+1+zr(2-r)\}\ge1,\quad z=c^2.
\]
Thus the six rational coefficients are the original coefficients by
uniqueness of the six stationary equations, including $n=0$.

The exact reflection is
\[
(u,v;W,X,Y,Z)\mapsto(1/v,1/u;Z,X,Y,W),\qquad
(m\ D)(C\ \alpha),
\]
and obeys, on the C2 omission face, the exact identity
\[
u^2v^2F_{c'}(1/v,1/u;Z,X,Y,W)=F_c(u,v;W,X,Y,Z).
\]
Both $n$ and $A$ are unchanged. The positive
source normalization by $W_c$ is legal. The factorization and endpoint
interpretation are homogeneous, so literal testing infinity is included.

\section{Exhaustive dispatch into the historical thirteen rows}

\subsection{C1: five retained rows and explicit other exits}

All entries in the following table include the exact C1 input above.
The notation $\neg\I$ means a positive unit or an unbounded element.
The row order is deterministic.
\begin{longtable}{p{.57\linewidth}p{.35\linewidth}}
\toprule Source condition & Destination \\\midrule
$k:\I$, $\zeta:\neg\I$ & C1.4, arbitrary positive $\eta$ \\
$k:\U$, $\eta,\zeta:\I$ & AS: arbitrary-state-ratio theorem \\
$k:\U$, at least one of $\eta,\zeta:\neg\I$ & C1.5 \\
$k,\zeta,\eta:\I$ & US: whole uniform-source exit; exact join below \\
$k,\zeta:\I$, $\eta:\Otype$ & C1.1 \\
$k,\zeta:\I$, $\eta:\U$, $e:\I$ & C1.3, every $0<r<1$ \\
$k,\zeta:\I$, $\eta,e,r:\U$ & BE: bounded-$e$ all-tests theorem \\
$k,\zeta,r:\I$, $\eta,e:\U$ & CA: centered-$A$, unit-$e$ theorem \\
$k,\zeta:\I$, $\eta,r:\U$, $e:\Otype$ & UE: unit-$\eta$, unbounded-$e$ theorem \\
$k,\zeta,r:\I$, $\eta:\U$, $e:\Otype$ & C1.2 \\
\bottomrule
\end{longtable}
To prove exhaustion, first split bounded $k$ into $\I/\U$. On the
$k:\I$ side split $\zeta$ into infinitesimal or not. On the remaining
$k,\zeta:\I$ side split $\eta$ into its three types. Only
$\eta:\U$ requires the $e$ trichotomy, and only $e:\U$ or
$\Otype$ requires the two possibilities for $r$. There is no
unbounded-$r$ case because $0<r<1$. The inequalities $a,t<2$ imply
$r\sim1$ when $\eta:\I$, but no extra rate condition is needed.

For AS, the original source formulas make $D$ and $\beta$ strictly
positive; a literal zero cannot be smuggled into its theorem. BE and CA
explicitly retain $\beta=0$; on their source shapes the formula for
$D$ has positive-unit numerator and prevents $D=0$. UE makes
$m/d$ and $\beta/d$ unbounded, so its stated assumptions also decide
that boundary. These are source consequences, not informal continuity
extensions.

The historical US entry has three exhaustive subentries. For $e:\I$
its exact theorem is \path{c1_uniform_all_tests.tex}. For $e:\Otype$
the original unbounded-$e$ $W$-unit/outer-test theorem and the restricted
contact--conductor completion provide the full normal restoration.
The new \path{us_unit/bounded_delta_all_tests.tex} proves the complete
source for bounded $\Delta=1/e$, hence covers both $e:\U$ and
$e:\Otype$. Its new grading assigns weight zero to both $\delta$
and testing $w$, while preserving strictly positive weights on every
scalar and first-normal correction coefficient. The $Z$-unit chart now
allows bounded $W$; the complementary $W$-unit, $Z$-infinitesimal chart
uses the retained-current identity with $\Delta Z$ infinitesimal.
The original outer-angle identities then give all projective tests.
This is the exact matching proof for US-$\U$ that the old summary
did not supply. The positive strict gaps and all original source signs
remain premises, and every subtracted packet is restored.

\subsection{C2: directed weak rows and strong orbits}

Weak means infinitesimal after division by $d$, while strong means
bounded positive reciprocal; strong coefficients may be unbounded.
Since $C,\alpha>0$, first split into both weak, exactly one weak, and
both strong. Both weak goes to the chart-free C2 radial theorem. Its
physical matching witnesses force all four matching weights large,
then the exact contact estimates give the radial-source entry. The
source argument and the final radial all-tests certificate are distinct
steps and are separately named in the ledger.

For exactly one weak coefficient reflect once, if necessary, so $C$
is weak and $\alpha$ strong. Do not quotient this directed list by
another reflection. Source identities exclude actual types II, IU, UI,
UU in that orientation. The all-tests single-angle corollary removes
UH. The remaining four historical directed rows are
\[
\mathrm{C2.W.IH},\quad\mathrm{C2.W.HI},\quad
\mathrm{C2.W.HU},\quad\mathrm{C2.W.HH}.
\]
Here H is an unbounded positive actual half-angle, not literal infinity;
the actual contact is always finite over $K_c$.

For both $C,\alpha$ strong, if $m,A,D$ are all strong then the five
consecutive events $(D,C,A,\alpha,m)$ give the five-event exit.
Otherwise retain $m$ weak, else $A$ weak, else $D$ weak, in that order
if a disjoint event subpartition is desired. The all-tests IU/UH and
UU corollary removes three ordered angle types. Reflection acts on
the other six by reversal and reciprocal type, giving precisely
\[
\mathrm{C2.S.II/HH},\quad\mathrm{C2.S.IH},\quad
\mathrm{C2.S.UI/HU},\quad\mathrm{C2.S.HI}.
\]
These are four orbits; IH and HI are each fixed and must not be
identified. Together with the five C1 rows this gives thirteen labels.
The finite angle calculation is included in the exact checker.

The directed HI and HU rows were subsequently proved empty, without
testing restrictions. The proof is substantive: with $q=u_c^{-2}:\I$
and bounded $y=v_c^2$, original $n\ge0$ forces $R=r/q$ bounded.
Otherwise the positive leading term in its exact numerator dominates
the negative terms, contradicting $n\ge0$. Then $H=v_cY_c\sim1$,
$\Delta$ is a positive bounded unit, and
$C=(1+R)(N_{\rm src}/q)/(2L\Delta)$ has bounded positive reciprocal.
This contradicts the chosen weak-$C$ orientation. Literal $n=0$ is
included in the contradiction. The ledger records the exact file.

The optional reciprocal-product exit requires all four products
$Cn,Cm,CA,CD$ to have bounded positive reciprocals in the weak-$C$
orientation. Merely two large adjacent weights do not suffice, and
$n=0$ never passes this exit. This exit is unnecessary for the finite
thirteen-row coverage implication; intersecting its complement later
preserves coverage.

\section{C1.4: exhaustive C14 subdivision, with missing exits restored}

The input is $k:\I$, $\zeta:\neg\I$, with all exact C1 constraints.
Put $L=x\zeta$ and, when $\eta:\Otype$, $\sigma=1/\eta$.
The following tree covers this whole source shape, before any extra
testing screens are applied.

\paragraph{$L:\Otype$.}
The strict original gap $Q=1-Lr>0$ forces $r:\I$.
Since $1-r=a\eta$ and $a<2$, $\eta$ cannot be infinitesimal.
For $\eta:\U$ the large-$L$, unit-$\eta$ all-tests theorem applies.
For $\eta:\Otype$, put $H=h/\sigma$. If $H:\Otype$, enter C14.1b.
If $H$ is bounded, split
$M_\beta=\beta\theta/\sigma^2$ into bounded or unbounded, after
$d=1$. Bounded $M_\beta$, including $\beta=0$, is the JR/JB all-tests
exit; unbounded $M_\beta$ is C14.1a. No $H=0$ case is introduced,
since $h>0$; $M_\beta=0$ is literal and allowed.

\paragraph{$L:\U$.}
The three types of $\eta$ give, respectively, the small-$\eta$
$L$-unit all-tests exit, the unit-$\eta$ $L$-unit all-tests exit, and
C14.2. The source formulas in the first two cited theorems, rather than
their titles, retain the original gap and zero-boundary hypotheses.

\paragraph{$L:\I$.}
Unbounded $\eta$ gives C14.4. For bounded $\eta$, unbounded $\zeta$
gives C14.3a. The only remaining case is $\zeta:\U$: infinitesimal
$\eta$ gives C14.3b. If both $\eta,\zeta:\U$, define the exact
positive bounded source gap
\[
j=\eta(1+xh)-1+r=\eta J.
\]
If $j:\U$, then $j_0=\eta+r-1$ is a positive unit and the RO theorem
applies. If $j:\I$, then $\eta+r\sim1$, hence $1-r:\U$.
For $r:\U$ use JO; for $r:\I$ use HO. There is no need to compare
$j$ with $x$, and no equality $j=0$ is permitted.

This proves that the six historical refined shapes
C14.1a, C14.1b, C14.2, C14.3a, C14.3b, C14.4, together with the exact
listed exits, cover C1.4. R19 admits all of C14.2 and C14.3a/b.
Their files are recorded as existing closures; this note does not claim
to have proved C14.1a/b or C14.4. A later continuation may replace
those open labels only with its actual proof bytes and scope.

\section{Weak HH: exact reflected entry and the ten disjoint rows}

Reflect the directed weak-$C$ HH source once to actual II. Now $C$
is strong and $\alpha$ weak. Put
\[
z=u_c^2,\quad y=v_c^2,\quad H=v_cY_c,\quad
L=v_c(u_cW_c-Y_c),\quad B=H+L,\quad r=-X_c/u_c.
\]
The exact original constraints include
\begin{gather*}
z,y:\I,\quad 0<H,r<1,\quad L>0,\quad j=1+zr-L/y>0,\\
G=B+z(H+r-1)>0,\quad B-zyr>0,\quad n\ge0.
\end{gather*}
Since $L<y(1+zr)$, $L$ is infinitesimal and $B$ is bounded.
The unchanged alpha formula is
\[
\alpha=\frac{(B+zr)M}{2(1-H)\delta},\qquad
M=(1+z)(B+zr)H+yG+yzj+zrL(1+y)>0,
\]
with the positive original denominator $\delta$ printed in the ledger.
If $B$ is infinitesimal, the exact estimates give $\alpha$ infinitesimal.
If $B$ is a unit, then $H$, $\delta$, and $M$ are units, so $\alpha$
is strong; also $CL=(1+zr)M/(2\delta)$ is a unit and $C$ is
unbounded. Thus the actual directed weak HH input has exactly
$B,z,y:\I$, not an assumed power-law scaling.

Define the weak-II variables
\[
a=z/B,\quad t=y/B,\quad h=H/B,\quad l=L/B=1-h,
\quad e=1-r,\quad g=G/B.
\]
All $h,l,r,e,g,j$ are strictly positive bounded elements; $a,t>0$ can
have all three types. Retain exactly
\[
g=1-ae+hz>0,\qquad j=1+zr-l/t>0,\qquad B-zyr>0,
\qquad n\ge0,
\]
as well as all original source formulas and $\delta_{\rm tot}:\I$.
The first identity gives
\[
ae<1+hz;\quad e:\U\Rightarrow a\text{ bounded};\quad
a:\Otype\Rightarrow e:\I;\quad a:\I\Rightarrow g\sim1;
\quad g:\I\Rightarrow ae\sim1. \tag{WF}
\]
The exact coefficient estimates WII1--WII4 prove that all six retained
event coefficients are bounded if and only if
$t,h,l,j,r,g$ are units. With $d=1$, that would make
$\delta_{\rm tot}$ a unit. Therefore the all-six-unit branch is
excluded on the present input. Also the two-weak source exit is exactly
$r,g,H/y,j:\I$; when $t:\U$ this is $r,g,h,j:\I$.

After those exits and the regular A-opening admissions, the ten ordered
rows are as follows. Every row includes all preceding weak-II equations
and original coefficients; a type table does not assert its realizability.
\begin{longtable}{c p{.78\linewidth}}
\toprule Row & Exact additional type condition \\\midrule
R1 & $t:\I$.\\
R2 & $t:\Otype$.\\
R3 & $t,h:\U$, $l:\I$.\\
R4 & $t,h,l:\U$, $j:\I$.\\
R5 & $t,h,l,j:\U$, $r:\I$.\\
R6 & $t,h,l,j,r:\U$, $g:\I$.\\
R7 & $t:\U$, $h,g:\I$, excluding $r,j$ both infinitesimal.\\
R8 & $t,g:\U$, $h,j:\I$.\\
R9 & $t,g,j:\U$, $h,e,a:\I$.\\
R10 & $t,g,j:\U$, $h,e:\I$, $a:\Otype$.\\
\bottomrule
\end{longtable}

Here is the proof of exhaustion, including every discarded type.
Split $t$ first. For $t:\U$, split $h$. If $h:\U$, first test
$l:\I$, then $j:\I$, then $r:\I$, then $g:\I$. These are R3--R6.
If none applies, all six variables $t,h,l,j,r,g$ are units, contradicting
the total-scale input by WII4. If $h:\I$, $l\sim1$. First test
$g:\I$: either $r,j:\I$ and the two-weak theorem applies, or this is
R7. For $g:\U$, test $j:\I$, giving R8. Only
$t,g,j:\U$, $h:\I$ remains.

On that final regular opening, if $e:\U$ then $a$ is bounded by
(WF). For $r:\U$, $a:\U$ gives WA and $a:\I$ gives WE; for
$r:\I$ the SL theorem applies for either bounded type of $a$.
If $e:\I$ then $r\sim1$; $a:\U$ gives WC, $a:\I$ gives R9,
and $a:\Otype$ gives R10. WA is joined to its printed remote theorem;
WE, SL with its projective join, and WC each have exactly the source
conditions just stated. No $h/a$ or $e/B$ comparison is inserted.

The table is disjoint by this order. The focused checker enumerates
all type assignments obeying the necessary consequences of (WF) and
$h+l=r+e=1$, and confirms that every surviving tuple has precisely
the displayed destination or one named exit. It does not assume that
every such combinatorial tuple is realized by a physical source.

R19 already admits R5, R8, R9, R10, and portions of R3, R4, R6, R7.
The supplied R19 frontier records the remaining exact intersections:
R3 and R4 have $e:\I$; R6 has $a:\Otype$, $e:\I$,
$\rho=n/e^2:\U$; R7 has $a:\Otype$, $e:\I$.
R1 and R2 retain their full original residual source conditions.
These are substitutions of existing theorem scopes after the coverage
argument, not evidence that a remaining domain is nonempty.

\section{Boundary, normalization, and testing rules}

\begin{enumerate}
\item C1 $D=0$ or $\beta=0$ is carried through the rational source
minor. If a later branch makes one impossible, its original coefficient
formula must prove that fact. The other zero is still kept. Both zero
simultaneously is outside the five/six-active coercive entry.
\item C2 $n=0$ is retained in global sign entry, the weak-II formula,
and every subsequent source inequality. It is never replaced by
$n>0$, nor used as a residue approximation to $n>0$.
\item A source shape parameter or gap is strictly positive unless its
definition is an original nonnegative coefficient. An infinitesimal
opening is not a literal zero. In particular $M,J,E,Q,r,1-r,h,l,g,j$
are not set to zero in an exact source formula.
\item $d=1$ and $C=2$ are legal positive pure-source normalizations;
they do not turn the earlier condition $\delta_{\rm tot}:\I$ into
a positive-total-scale source. Source-only coefficient weights and
testing-only normalization factors have square roots in their respective
fields and restore every certificate.
\item The source angles classify actual contacts. Independent testing
angles remain arbitrary, with signs, zeros, and infinity treated in their
homogeneous charts. A finite-chart proof alone is not an all-tests
statement. The ledger distinguishes local theorems from their exact joins.
\item The optional rational-curve theorem requires coefficients in a
single $\R(T)$, a specified physical real interval, and nonvanishing
denominators there. It is not an arbitrary multiparameter exit.
The C2-to-C1 overlap is physical only for $n=0$: the overlap identity
gives a negative C1 beta coefficient for $n>0$.
\end{enumerate}

\section{Conclusion of this audit, with the remaining entries explicit}

The omission routing, positive-chart maps, C1/C2 source-type disjunction,
C14 refinement, and weak R1--R10 refinement are now explicit. Their
exhaustion is proved on the algebraic source input with the stated type
conditions, independently of nonemptiness and independently of optional
historical exclusions. The exact one/two-hard and four-matching work
completes the requested U1.A6 technical audit on its stated domain.
The separate bounded-$\Delta$ proof repairs the US-$\U$ terminal
statement mismatch without relying on a historical global assertion.

The precise remaining audit entries are:
\begin{enumerate}
\item U1.A5: the scale split and source partition were already audited
in R18. The ordinary and finite-pitchfork subbundle with
$d,C,A,D,\alpha$ units is now repaired using the recovered full-parent
finite-edge theorem, verified exact source entry, and a complete
projective join. The separate note displays all four corner quadratic
forms and restores every retained packet. Its exact symmetry images
are included only under their stated unit hypotheses. The separate C2
weak-$\alpha$ note proves the subdomain with all six of
$d,n,m,C,A,D$ units and its exact reflection, by showing every
remaining singular zero-corner jet fits the recovered finite-ratio
theorem. No bound on the actual tested ratio is assumed. Remaining
weak-alpha shapes and their stabilizer images, the C1 projective
pitchfork, and reciprocal all-tests terminals remain to be accepted.
\item U1.A7: the ledger now names essential and optional exits, but
their complete all-tests proofs are not all freshly reaudited here.
The former US-$\U$ provenance issue is repaired by the new full-parent
theorem for $k,\eta,\zeta:\I$, bounded $1/e$, all original gaps and
signs, and all independent tests. Its review is recorded separately
from the imported terminal bundle.
\item U1.A9: parts (i)--(iv) have the explicit coverage construction
above. Its terminal-substitution step remains conditional on the exact
registry inputs and on the open local rows in the
current frontier. Do not change a historical import to ``newly proved''
because its filename appears in the registry.
\end{enumerate}
U1 remains open; the global representability bridge remains blocked by
U1. No separate uniform bound on the number of local squares has been
introduced. This note is a same-system mathematical audit with narrow
exact symbolic corroboration, not external review or proof-assistant
verification.
